\documentclass[11pt,a4paper]{article}
\usepackage[T1]{fontenc}
\usepackage{lmodern,amsmath,amssymb,booktabs,microtype}
\usepackage[margin=25mm]{geometry}
\usepackage[colorlinks=true,urlcolor=blue,linkcolor=blue]{hyperref}
\newcommand{\E}{\mathbb E}
\newcommand{\K}{\mathcal K}
\newcommand{\R}{\mathcal R}
\title{Real principal probes in the dS$_3$ eikonal model:\\lateral kernels and Gaussian spectral moments}
\author{Alan Day \and Chengheng\thanks{AI research assistant (OpenAI).}}
\date{Working research note --- 31 October 2026}
\begin{document}\maketitle
\begin{abstract}
We retain the Bessel wavefunctions of real principal-series probes in
the existing observer eikonal prescription. Opposite momentum rotations
fix the lateral continuation. A four-rapidity representation permits
four independent real Gaussian time smearings without discarding the
complex phase. The averaged momentum moments have a positive Mellin
spectral representation, even though the unaveraged integrands change
sign. This fixes the signs of every fixed-order coefficient, but not
the nonperturbative imaginary part. Independent momentum integrals and
scrambled Sobol calculations test the formulas. These are results within
the specified approximation, not a finite-$G$ Hilbert-space construction
or a completed holographic dictionary.
\end{abstract}

\section{Scope and conventions}
Use radius $\ell=1$, observer mass $m>0$, $\gamma=4Gm\in[0,1)$,
and $\alpha=2\pi(1-\gamma)$. The two distinguishable scalar probes have
$\nu,\nu'>0$ real and mass squares $1+\nu^2,1+\nu'^2$.
The state/preparation is the fixed-average-observer-mass mixed
representation of CSTY, with its prescribed $12/14$ contours
\cite{csty}. It is not already identified with the preceding
$G=0$ quantized-observer Fock vacuum. We sum the two eikonal sectors
and subtract the duplicated disconnected term, as in CSTY (2.25).
The original wavefunctions, shock, and sector construction are prior
results. Our exercise applies standard integral identities to their
four-time Gaussian smearing; no priority claim is made.

Let $d_i$ be independent centered Gaussians of variance $\sigma^2$,
with fixed $\sigma>0$, and set
\[
 (\theta_1,\theta_2,\theta_3,\theta_4)
 =(-\pi/2+i(t+d_1),\,id_2,\,\pi/2+i(t+d_3),\,\pi+id_4).
\]
Write $X=(d_1-d_3)/2$, $Y=(d_2-d_4)/2$,
$Z=(d_1+d_3-d_2-d_4)/2$ and $T=t+Z$.
These are independent, of variances $\sigma^2/2,\sigma^2/2,\sigma^2$.
Angular brackets mean $\pi^{-1}\int_0^\pi d\phi$, with
\[
 h(\phi)=\frac{H_\gamma(\phi)}m,
 \qquad H_\gamma(\phi)=-\frac{\cos[(1-\gamma)(\pi-\phi)]}
 {2\operatorname{sinc}(\pi\gamma)},\quad \operatorname{sinc}x=\frac{\sin x}x.
\]

\newpage
\section{Momentum rotations and the lateral kernel}
Define the unnormalized positive-coupling integral
\begin{align}
 S(\lambda;X,Y)&=\frac{1}{(\pi\alpha)^2}
 \int_0^\infty x\,dx\,y\,dy\,P_\nu(x,Y)P_{\nu'}(y,X)e^{-\lambda xy},\label{eq:S}\\
 P_\nu(x,Y)&=K_{i\nu}(xe^Y)K_{i\nu}(xe^{-Y}),\qquad \Re\lambda\geq0.\nonumber
\end{align}
This integral diverges on the negative real $\lambda$ axis and cannot
be used there without continuation. Here is a direct check of that
continuation from CSTY (2.32). Put $s_{13}=(d_1+d_3)/2$,
$s_{24}=(d_2+d_4)/2$, and rotate its negative momenta oppositely:
\[
 p_+=-e^{i\pi/4-s_{24}}x,\qquad
 p_-=-e^{-i\pi/4+t+s_{13}}y,\qquad z=e^{-i\pi/4}.
\]
Their product stays real. The source integral becomes
\[
 -\frac{1}{(\pi\alpha)^2}\int_0^\infty x\,dx\,y\,dy\,
 P_\nu(zx,Y)P_{\nu'}(zy,X)e^{i\lambda xy},\quad \lambda=he^T.
\]
The minus sign is fixed by the source prefactors and Jacobians.
For real $\lambda$ this integral is absolutely convergent. Using the
standard representation \cite{dlmfK}
\[
 K_{i\nu}(w)=\int_0^\infty e^{-w\cosh r}\cos(\nu r)\,dr,
 \quad \Re w>0,
\]
introduce $r_j\geq0$ and
\[
 A=e^Y\cosh r_1+e^{-Y}\cosh r_2,\quad
 B=e^X\cosh r_3+e^{-X}\cosh r_4.
\]
\[
 W=\cos\nu r_1\cos\nu r_2\cos\nu'r_3\cos\nu'r_4.
\]
Integrating $y$ first and setting $v=zAx$ leaves
$(A^2B^2)^{-1}\int v e^{-v}(1+cv)^{-2}dv$, $c=\lambda/(AB)$,
on the ray $\arg v=-\pi/4$. Rotation to the positive axis passes
\emph{below} the positive pole for $c<0$. Thus
\begin{align}
 S_\eta(\lambda;X,Y)&=\frac{1}{(\pi\alpha)^2}
 \int_{\mathbb R_+^4}\frac{W\,d^4r}{A^2B^2}\K_\eta\left(\frac\lambda{AB}\right),\label{eq:rapid}\\
 J(c)&=\frac{e^{1/c}E_1(1/c)}c,\qquad
 \K(c)=-J'(c)=\frac{(1+c)J(c)-1}{c^2},\quad \K(0)=1.\nonumber
\end{align}
The above rotated source is $\eta=+1$; the other boundary value is its
conjugate. For $c=-1/a$, $a>0$,
\[
 J_\eta(-1/a)=ae^{-a}\operatorname{Ei}(a)-\eta i\pi ae^{-a},\qquad
 \Im\K_\eta(-1/a)=-\eta\pi a^2(a-1)e^{-a}.
\]
Both this residue and $W$ change sign. The conformal sign proof does
not extend to real principal probes by replacing their masses.

\newpage
\section{A positive spectrum for the smeared moments}
For every integer $n\geq0$ define
\[
 I_{n,\nu}(X)=\int_0^\infty x^{n+1}P_\nu(x,X)\,dx,
 \qquad B_{n,\nu}(\sigma)=\E_X I_{n,\nu}(X).
\]
Set $f_n(r)=e^{(n+2)r/2}K_{i\nu}(e^r)$. This is a real $L^2$
function, and changing to logarithmic momentum gives
\[
 I_{n,\nu}(X)=\int_{\mathbb R}f_n(r+X)f_n(r-X)\,dr.
\]
With Fourier transform $\widehat f(\omega)=\int e^{-i\omega r}f(r)dr$,
the Mellin identity \cite{dlmfM} yields
\[
 M_{n,\nu}(\omega)=2^{n/2-1-i\omega}
 \Gamma\left(\frac{n+2}{4}+\frac{i(\nu-\omega)}2\right)
 \Gamma\left(\frac{n+2}{4}-\frac{i(\nu+\omega)}2\right).
\]
Its Mellin convergence condition holds since $(n+2)/2>0$ and
$\Re(i\nu)=0$. Plancherel and
$\E e^{2i\omega X}=e^{-\sigma^2\omega^2}$ prove
\begin{equation}
 \boxed{\quad B_{n,\nu}(\sigma)=\frac1{2\pi}\int_{\mathbb R}
 e^{-\sigma^2\omega^2}|M_{n,\nu}(\omega)|^2d\omega>0.\quad}\label{eq:B}
\end{equation}
Fubini is justified by the integrable nonnegative Fourier density.
This positivity is a statement about these averaged moments, not
positivity of $P_\nu$ or of a complete gravitational Hilbert space.
The variable $\omega$ is Fourier conjugate to logarithmic momentum;
no identification with modular energy is implied.

Two normalization checks are
\begin{align*}
 I_{0,\nu}(X)&=\frac{\pi\sin(2\nu X)}{2\sinh(\pi\nu)\sinh(2X)},
 &I_{0,\nu}(0)&=\frac{\pi\nu}{2\sinh\pi\nu},\\
 I_{n,\nu}(0)&=\frac{2^{n-1}\Gamma((n+2)/2)^2
 |\Gamma((n+2)/2+i\nu)|^2}{\Gamma(n+2)}.
\end{align*}
In particular, $I_{0,\nu}(X)$ can be negative, whereas its Gaussian
average is strictly positive. The disconnected normalization is
\begin{equation}
 D_\sigma=\frac{B_{0,\nu}(\sigma)B_{0,\nu'}(\sigma)}{(\pi\alpha)^2}.
\end{equation}
The complete specified two-sector model is
\begin{equation}
 \R_{\eta,\sigma}(t)=\frac{\E\langle
 S_\eta(he^{t+Z};X,Y)+S_{-\eta}(he^{-t-Z};X,Y)
 \rangle-D_\sigma}{D_\sigma}.\label{eq:R}
\end{equation}
There is no pointwise division by the sign-changing two-point
functions. The model obeys $\R_{-\eta,\sigma}=\overline{\R_{\eta,\sigma}}$
and $\R_{\eta,\sigma}(-t)=\overline{\R_{\eta,\sigma}(t)}$.

\newpage
\section{Fixed orders, signs, and limits}
The fixed-order expansion of the elementary kernel is
\[
 \K_\eta(c)\sim\sum_{n\geq0}(-1)^n(n+1)(n+1)!c^n.
\]
At fixed $\gamma,\sigma,t,\nu,\nu'$, therefore,
\begin{equation}
 \R_{\eta,\sigma}\sim1+\sum_{n\geq1}
 \frac{(-1)^n\langle h^n\rangle}{n!}
 e^{n^2\sigma^2/2}
 \frac{B_{n,\nu}B_{n,\nu'}}{B_{0,\nu}B_{0,\nu'}}
 (e^{nt}+e^{-nt}).\label{eq:series}
\end{equation}
The moment expansion before smearing is already present in CSTY
(2.37)--(2.38); related momentum moments occur in MNX (5.23)--(5.28)
\cite{mnx}. Equation \eqref{eq:B} is what fixes the sign after the
independent Gaussian averages.

For $0<\gamma<1$, every even shock moment is positive and every odd
shock moment is negative. To see the latter, put $b=(1-\gamma)\pi$.
The integral $\int_0^b\cos^{2k+1}v\,dv$ is positive for $0<b<\pi$:
for $b\leq\pi/2$ this is direct; otherwise use the vanishing full
integral on $[0,\pi]$ and the negative omitted piece. Hence every
coefficient in \eqref{eq:series} is positive. At $\gamma=0$ the odd
ones vanish and the even ones stay positive. This is not convergence
of the series, monotonicity of the full answer, or a statement about
the sign of its imaginary part.

For completeness, the termwise operation is valid as a fixed-order
asymptotic expansion. The real lateral kernel is continuous at zero,
bounded at finite $c$, and tends to zero as $|c|\to\infty$.
Its fixed-order small-$c$ remainder has the usual $E_1/\operatorname{Ei}$
bound; the exponentially small cut term is flat there. Thus, for
each $N$, a finite constant $d_N$ exists with
\[
 \left|\K_\eta(c)-\sum_{n=0}^{N}(-1)^n(n+1)(n+1)!c^n\right|
 \leq d_N|c|^{N+1}\qquad(c\in\mathbb R).
\]
On $|c|$ bounded away from zero this follows directly by boundedness
and polynomial growth. In \eqref{eq:rapid}, use $|\cos|\leq1$ and
the autocorrelation bound $|I_{N+1,0}(X)|\leq I_{N+1,0}(0)$.
The normalized remainder in \eqref{eq:series} is at most
\[
 \frac{2d_N\|h\|_\infty^{N+1}e^{(N+1)^2\sigma^2/2}
 \cosh((N+1)t)}{B_{0,\nu}B_{0,\nu'}}
 \left[\frac{I_{N+1,0}(0)}{(N+2)!}\right]^2.
\]
No useful uniformity as $N\to\infty$ is asserted.

As a check, at the separate conformal value $\nu=i/2$,
$I_{n,i/2}(X)=\pi n!/[2^{n+2}\cosh^{n+1}X]$. This recovers the earlier
conformal smearing factors, without analytically continuing a
Hilbert-space theorem. At fixed $t$, $m\to\infty$ gives $\R\to1$.
At $t=\log m+u$ the forward sector instead keeps the argument
$H_\gamma e^{u+Z}$; the reverse sector cancels the subtraction.

\newpage
\section{Absolute integrability and numerical evidence}
Let $C_K=\sup_{c\in\mathbb R,\eta=\pm1}|\K_\eta(c)|<\infty$.
We do not assign a numerically certified value to this constant.
The exact nonoscillatory majorant is
\[
 \int_0^\infty\!\int_0^\infty
 \frac{dr\,ds}{(e^X\cosh r+e^{-X}\cosh s)^2}
 =\frac{X}{\sinh2X}\leq\frac12.
\]
Consequently $|S_\eta|\leq C_K/[4(\pi\alpha)^2]$, uniformly in
real $X,Y,\lambda$. Gaussian smearing and the two limits above follow
by domination. A cutoff at $L$ standard deviations in each of
$X,Y,Z$ omits normalized absolute mass at most
\[
 \frac{3(2C_K+1)}{4B_{0,\nu}B_{0,\nu'}}\operatorname{erfc}(L/\sqrt2).
\]
For a rapidity cutoff $r_j\leq R$, use
$\cosh r+\cosh s\geq(e^r+e^s)/2$ and split $s<r$, $s>r$.
The pair union tail is at most $4(R+1)e^{-2R}e^{2|X|}$.
Since $\E e^{2|X|}=e^{\sigma^2}\operatorname{erfc}(-\sigma)$,
the full normalized rapidity tail is bounded by
\[
 \frac{4(2C_K+1)(R+1)e^{-2R}e^{\sigma^2}
 \operatorname{erfc}(-\sigma)}{B_{0,\nu}B_{0,\nu'}}.
\]
These continuum bounds do not certify discretization, arithmetic,
or terms omitted by the gravity approximation.

We use $\nu=.6$, $\nu'=.9$, $m=1000$ and $t=\log m+u$.
The denominator is integrated separately at 50 decimal digits by
\eqref{eq:B}. Four rapidities are sampled with exponential importance
density and three normal variables with inverse Gaussian CDFs; the
angle is uniform. Signed cosine weights and both lateral kernels
are retained. Eight independent scrambles at $2^{18}$ and $2^{20}$
points are followed by eight new scrambles at $2^{21}$ points with
importance rate $.7$ instead of $1$. Raw and control-variate estimates
are both saved. The control variate replaces only the sampled
disconnected term by its independently known integral.

\begin{center}\begin{tabular}{ccc rr rr}
\toprule
$\sigma$&$\gamma$&$u$&$\Re\R_+$&$\Im\R_+$&SE$_{\Re}$&SE$_{\Im}$\\
\midrule
.4&0&0&1.10266&$-.04401$&$.000024$&$.000032$\\
.4&.25&0&1.24688&$-.08601$&$.000021$&$.000052$\\
.8&.25&0&1.19389&$-.19901$&$.000055$&$.000101$\\
.8&.25&2&.30011&$-.81574$&$.000239$&$.000280$\\
\bottomrule
\end{tabular}\end{center}
SE denotes the empirical standard error across eight scrambles,
not a rigorous error bar. Changing rate and seeds shifts the listed
components by at most $2.68\times10^{-4}$; refinement is not monotone.
All 60 initial diagnostics and 14 rotation/refinement checks pass
their stated numerical tolerances. These include two positive-coupling
full Bessel integrals, three oppositely rotated complex momentum
integrals, and conformal normalization controls. The largest smeared
moment relative comparison error is $1.88\times10^{-11}$.

\newpage
\section{Physical interpretation and remaining work}
The nonzero imaginary parts in the table are numerical evidence for
the specified contour model. There is no principal-series analogue
of our conformal all-time sign theorem here: already at $X=Y=0$,
the unnormalized rapidity integral changes imaginary sign between
$\lambda=-1$ and $-10$ in the saved deterministic quadratures.
A separate full scan at $\sigma=.4,\gamma=.25,u=4$ gives
$\Im\R_+=.02768$ with empirical SE $.000056$, while $u=0$ is negative.
At $u=6$ the raw estimate is $.001038-.000388i$, with component
SEs $.000037,.000031$ and absolute-integrand average $.1941$;
cancellation is substantial. At $u=8$ the imaginary estimate
$.000034$ has SE $.000024$: its sign is unresolved. All seeds are
saved, and none of these empirical observations is a universal sign law.
For a Hermitian equal-spacing thermal trace, whenever the requisite
products are trace class, cyclicity and adjunction require a real
answer. The present computation preserves rather than resolves the
known contour issue. It does not exclude omitted gravitational
contributions, different states, or a different physical prescription.

Recent adjacent work does not currently supply the missing map.
MNX distinguishes a graviton-mass-regulated $L=1$ double scaling from
an unregulated leading tree result \cite{mnx}. That regulator is not
our finite observer recoil. The horizon-edge EFT study \cite{edge}
uses a $D\geq4$ test-observer setting and leaves quantization and the
coupling to a physical observer open. Its corner choices cannot be
silently substituted for a $dS_3$ massive-observer state or contour.
The literature reading record states which sections were inspected;
we do not claim complete reproduction of either paper.

The next physical task is to compare the actual state, normalization,
insertion and gluing data with a microscopic candidate. Agreement
of these eikonal moments alone is insufficient. In particular,
the finite-parameter insertion normalization and Gram inverse,
the correspondence with the regulated $G=0$ Fock vacuum, and
control of omitted gravity terms remain unresolved.

\paragraph{Reproduction.}
The code and immutable raw runs live under \path{code/} and
\path{results/stage4/principal-eikonal/}. Run the one-shot drivers
only in a clean output directory; they refuse to overwrite evidence.
\path{study.json} records dependency versions and input hashes.
\path{diagnostics.json} and \path{refinement.json} specify the exact
checks and tolerances. The stage evidence binds this source, the
compiled PDF, the independent controls and the saved outputs.
Prior scientific stages are retained byte for byte.

\begin{thebibliography}{9}
\bibitem{csty} Y. Chen, D. Stanford, H. Tang and Z. Yang,
\emph{Negative shocks versus static patch holography},
\href{https://arxiv.org/abs/2607.14042v1}{arXiv:2607.14042v1},
especially (2.25), (2.32)--(2.40) and Appendix G.
\bibitem{dlmfK} NIST Digital Library of Mathematical Functions,
\href{https://dlmf.nist.gov/10.32.E9}{10.32.9} (Bessel integral).
\bibitem{dlmfM} NIST Digital Library of Mathematical Functions,
\href{https://dlmf.nist.gov/10.43.E19}{10.43.19} (Mellin transform).
\bibitem{mnx} A. Milekhin, V. Narovlansky and J. Xu,
\emph{Out-of-time-ordered Correlators in de Sitter Revisited},
\href{https://arxiv.org/abs/2607.13137v1}{arXiv:2607.13137v1},
Sections 3--6.
\bibitem{edge} J. Bonifacio, K. Hinterbichler, Y. T. A. Law and V. Lochab,
\emph{De Sitter Observers and an Effective Field Theory of Horizon Edge Modes},
\href{https://arxiv.org/abs/2610.28719v1}{arXiv:2610.28719v1},
Sections 2--6 and Appendix A; PDF used for the truncated HTML portion.
\end{thebibliography}
\end{document}
